\documentclass[a4paper, 12pt]{article}
\usepackage[left=2cm,right=2cm,top=2cm,bottom=2cm]{geometry}
\usepackage{wrapfig} \usepackage{graphicx} \usepackage{mathtext}
\usepackage{amsmath} \usepackage{siunitx} % Required for alignment
\usepackage{subfigure} \usepackage{multirow} \usepackage{rotating}
\usepackage[T1,T2A]{fontenc} \usepackage[russian]{babel} \usepackage{caption}
\usepackage{hyperref}
\usepackage[]{float}

\hypersetup{
    colorlinks=true,
    linkcolor=red,
    urlcolor=magenta,
    }
\graphicspath{{pictures/}}


\title{\begin{center}Лабораторная работа №2.2.6\end{center} Определение энергия
активации по температурной зависимости вязкости жидкости} \author{Балдин Виктор
Б01-303} \date{\today}

\begin{document}
\maketitle

\textbf{Цель работы:} 1) измерение скорости падения шариков при разной
температуре жидкости; 2) вычисление вязкости жидкости по закону Стокса и расчет
энергии активации.

\textbf{В работе используются:} стеклянный цилиндр с исследуемой жидкостью
(глицерин); термостат; секундомер; горизонтальный компаратор; микроскоп; мелкие
шарики (диаметром 1 -- 2 мм).

\section{Теоретическая часть} \subsection{Энергия активации} Для того чтобы
перейти в новое состояние, молекула жидкости должна преодолеть участки с большой
потенциальной энергией, превышающей среднюю тепловую энергию молекул. Для этого
тепловая энергия молекул должна — вследствие флуктуации — увеличиться на
некоторую величину $W$ , называемую энергией активации. Температурная
зависимость вязкости жидкости при достаточно грубых предположениях можно описать
формулой \begin{equation} \label{activation_energy:1} \eta \sim A e^{W/kT}
\end{equation}

Из формулы (\ref{activation_energy:1}) следует, что существует линейная
зависимость между величинами $ln\eta$ и $1/T$, и энергию активации можно найти
по формуле

\begin{equation} \label{activation_energy:2} W = k \frac{d(\ln\eta)}{d(1/T)}
\end{equation}

\subsection{Измерение вязкости} По формуле Стокса, если шарик радиусом $r$ и со
скоростью $v$ движется в среде с вязкостью $\eta$, и при этом не наблюдается
турбулентных явлении, тормозящую силу можно найти по формуле (\ref{stokes})

\begin{equation}\label{stokes} F = 6\pi\eta rv \end{equation}


Для измерения вязкости жидкости рассмотрим свободное падение шарика в жидкости.
При медленных скоростях на шарик действуют силы Архимеда и Стокса, выражения для
которых мы знаем. Отсюда находим выражения для установившейся скорости шарика и
вязкости жидкости

\begin{eqnarray} v_{уст} =\frac{2}{9}gr^2\frac{\rho - \rho_{\text{ж}}}{\eta}\label{v_ust},\\
    \eta    =\frac{2}{9}gr^2\frac{\rho - \rho_{\text{ж}}}{v_{\text{уст}}}\label{eta}
\end{eqnarray}

Как видим, измерив установившуюся скорость шарика и параметры системы можно
получить вязкость по формуле (\ref{eta}).

\subsection{Экспериментальная установка} Для измерений используется стеклянный
цилиндрический сосуд B, наполненный исследуемой жидкостью (глицерин). Диаметр
сосуда $\approx 3$ см, длина $\approx 25$ см. На стенках сосуда нанесены две
метки на некотором расстоянии друг от друга. Верхняя метка должна располагаться
ниже уровня жидкости с таким расчетом, чтобы скорость шарика к моменту
прохождения этой метки успевала установиться. Измеряя расстояние между метками,
b время падения определяют установившуюся скорость шарика $v_{уст}$. Сам сосуд B
помещен в рубашку D, омываемую водой из термостата. При работающем термостате
температура воды в рубашке D, а потому и температура жидкости 12 равна
температуре воды в термостате.  Схема прибора (в разрезе) показана на
рис.~\ref{stand}.  \begin{figure}[h]
\center{\includegraphics[scale=0.5]{stand.png}} \caption{Установка для
определения коэффициента вязкости жидкости.} \label{stand} \end{figure}


\section{Ход работы} \subsection{Подготовительные работы} Для начала отбираем
примерно 25 шариков, и измеряем их диаметры. Диаметры измеряем в трех случайных
направлениях и усредняем. Это делается по той причине, что некоторые шарики (в
частности металлические) имеют неидеальную геометрию. Данные измерении приведены
в таблице \ref{diams}. В таблице шарики с номерами вида $s\#$ стеклянные, а вида
$m\#$ металлические. Погрешности измерении диаметров $\sigma_d = 0.02\ \text{мм}$.
Плотности шариков в эксперименте

\begin{align*} \rho_{\text{стекло}} & =2.5\ \text{г}/\text{см}^3 \\ \rho_{\text{металл}} & =7.8\ \text{г}/\text{см}^3
\end{align*}


\begin{table}[h!]
    \vspace{5pt}
    \begin{center}
        \subtable
        {
            \begin{tabular}{|l|rrr|r|}
                \hline
                №   & $d_1$ & $d_2$ & $d_3$ & $\langle d \rangle$ \\
                    & мм    & мм    & мм    & мм                  \\
                \hline
                s1  & 2.10  & 2.04  & 2.10  & 2.08                \\
                s2  & 2.08  & 2.06  & 2.06  & 2.07                \\
                s3  & 2.10  & 2.10  & 2.10  & 2.10                \\
                s4  & 2.06  & 2.08  & 2.08  & 2.07                \\
                s5  & 2.06  & 2.08  & 2.06  & 2.07                \\
                s6  & 2.08  & 2.08  & 2.08  & 2.08                \\
                s7  & 2.04  & 2.06  & 2.08  & 2.06                \\
                s8  & 2.04  & 2.10  & 2.08  & 2.07                \\
                s9  & 2.10  & 2.08  & 2.06  & 2.08                \\
                s10 & 2.08  & 2.08  & 2.08  & 2.08                \\
                s11 & 2.10  & 2.10  & 2.10  & 2.10                \\
                m1  & 0.66  & 0.70  & 0.70  & 0.69                \\
                m2  & 0.68  & 0.66  & 0.68  & 0.67                \\
                \hline
            \end{tabular}
        }
        \subtable
        {
            \begin{tabular}{|l|rrr|r|}
                \hline
                №   & $d_1$ & $d_2$ & $d_3$ & $\langle d \rangle$ \\
                    & мм    & мм    & мм    & мм                  \\
                \hline
                m3  & 0.84  & 0.84  & 0.84  & 0.84                \\
                m4  & 0.82  & 0.82  & 0.82  & 0.82                \\
                m5  & 0.76  & 0.76  & 0.76  & 0.76                \\
                m6  & 0.82  & 0.88  & 0.72  & 0.81                \\
                m7  & 0.92  & 0.88  & 0.94  & 0.91                \\
                m8  & 0.88  & 0.92  & 0.88  & 0.89                \\
                m9  & 0.90  & 0.90  & 0.88  & 0.89                \\
                m10 & 0.88  & 0.88  & 0.86  & 0.87                \\
                m11 & 0.96  & 0.92  & 0.92  & 0.93                \\
                m12 & 1.00  & 1.04  & 0.88  & 0.97                \\
                m13 & 0.88  & 0.90  & 0.88  & 0.89                \\
                m14 & 0.94  & 0.98  & 1.00  & 0.97                \\
                m15 & 0.90  & 0.90  & 0.90  & 0.90                \\
                \hline
            \end{tabular}
        }

        \caption{Измеренные диаметры шариков.}
        \label{diams}
    \end{center}
\end{table}

Измеряем длины частей цилиндра установки (см. рис. \ref{ustanovka})
\begin{equation*} l_1=l_2=(10.2\pm0.1)\ \text{см} \end{equation*}

\subsection{Измерение установившихся скоростей} Мы знаем путь, который проходит
шарик от одной отметки цилиндра к другой. Осталось измерить время прохождения
между этими отметками для получения скорости. В данной работе время падения
определяется секундомером. Получаем следующие данные (см.
таблицу \ref{times_in_frames}). Как видим $t_1$ и $t_2$ всегда
близки. Отсюда можно предположить что на рассматриваемых участках скорость не
меняется. В дальнейшем будем считать это предположение правдивым, которое в
дальнейшем подтверждается малостью времени и пути релаксации.

\begin{table}[H] \vspace{5pt} \begin{center} \subtable {
\begin{tabular}{|l|r|r|r|} \hline {№} & {$T,\ \text{\textdegree C}$} & {$t_1,\ \text{c}$} & {$t_2,\ \text{c}$} \\
\hline s1  & 30.86           & 15.3     & 15.1     \\ s2  & 30.86           & 15.2
& 15.1     \\ s3  & 41.53           & 8.5     & 8.5     \\ s4  & 41.50
& 8.5     & 8.4     \\ s5  & 50.93           & 4.9     & 4.9     \\ s6  & 51.05
& 4.6     & 4.7     \\ s7  & 62.00           & 2.8      & 2.9      \\ s8  & 61.97
& 2.7      & 2.9      \\ s9  & 61.85           & 2.6      & 2.6      \\ s10 & 68.40
& 2.0      & 2.0      \\ s11 & 68.47           & 2.0      & 2.0      \\ \hline
\end{tabular} } \subtable { \begin{tabular}{|l|r|r|r|} \hline {№} & {$T,\ \text{\textdegree C}
$} & {$t_1,\ \text{c}$} & {$t_2,\ \text{c}$} \\ \hline m3  & 30.84           & 16.8     & 16.8     \\
m4  & 30.85           & 17.4     & 17.4     \\ m5  & 41.51           & 10.7     &
10.4     \\ m6  & 41.53           & 10.0     & 10.0     \\ m7  & 51.09           &
4.4     & 4.5     \\ m8  & 51.05           & 4.5     & 4.4     \\ m10 & 61.53
& 2.8       & 2.7      \\ m11 & 61.16           & 2.7      & 2.7      \\ m12 & 61.16
& 2.4      & 2.4      \\ m14 & 67.81           & 2.0      & 2.0      \\ m15 & 67.73
& 2.0      & 2.0      \\ \hline \end{tabular} }

\caption{Измеренные времена падения шариков в кадрах.} \label{times_in_frames}
\end{center} \end{table} Для каждого измерения считаем $v$, $\eta$, $Re$,
$\tau$, $S$ где $v$ это скорость шарика на участке 1+2, $\tau$ это время
релаксации (см. формулу \ref{relaxation_time}), а $S=v\tau$ это путь релаксации.
\begin{equation}\label{relaxation_time} \tau = \frac{2r^2\rho}{9\eta}
\end{equation}

Плотность жидкости берем из графика \ref{density} \begin{figure}[h]
\center{\includegraphics[width=0.7\linewidth]{density}} \caption{Плотность
глицерина при различных температурах.} \label{density} \end{figure} \paragraph{}
Данные всех расчетов приведены в таблице \ref{data}

\begin{table}
    \begin{center}
        \begin{tabular}{|l|l|rr|rr|lll|}
            \hline
            {№} & {$T,\ \text{\textdegree C}$} & {$v,\ \text{см}/\text{с}$} & {$\Delta v, \text{см}/\text{с}$} & {$\eta, \text{мПа} \cdot \text{с}$} & {$\Delta\eta,\ \text{мПа}\cdot\text{с}$} & {Re} & {$\tau,\ \text{мс}$} & {$S,\ \mu\text{м}$} \\
            \hline
            s1  & 30.9          & 0.67        & 0.01              & 441            & 21                   & 0.02 & 0.7         & 4.7        \\
            s2  & 30.9           & 0.67        & 0.01              & 435            & 8                    & 0.02 & 0.7         & 4.7        \\
            s3  & 41.5           & 1.20        & 0.01              & 251            & 3                    & 0.06 & 1.2         & 14.4         \\
            s4  & 41.5           & 1.21        & 0.01              & 240            & 5                    & 0.07 & 1.2         & 14.4        \\
            s5  & 50.9           & 2.05        & 0.02              & 143            & 3                    & 0.18 & 2.1         & 43.1         \\
            s6  & 51.0           & 2.13        & 0.02              & 139            & 2                    & 0.20 & 2.2         & 46.9        \\
            s7  & 62.0           & 3.52        & 0.05              & 83             & 3                    & 0.55 & 3.5         & 123.2        \\
            s8  & 62.0           & 3.61        & 0.05              & 81             & 4                    & 0.57 & 3.6         & 130.0         \\
            s9  & 61.9           & 3.79        & 0.05               & 78             & 2                    & 0.63 & 3.8         & 144.0         \\
            s10 & 68.4           & 5.01        & 0.08              & 59             & 1                    & 1.09 & 5.0        & 250.1         \\
            s11 & 68.5           & 5.05        & 0.08              & 60             & 1                    & 1.10 & 5.1         & 257.6         \\\hline
            m3  & 30.8           & 0.60        & 0.01              & 420            & 4                    & 0.01 & 0.1         & 0.6      \\
            m4  & 30.9           & 0.58        & 0.01              & 414            & 4                    & 0.01 & 0.1         & 0.6       \\
            m5  & 41.5           & 0.96        & 0.01               & 215            & 2                    & 0.02 & 0.2         & 1.9        \\
            m6  & 41.5           & 1.01        & 0.01               & 233            & 66                   & 0.02 & 0.2        & 1.9        \\
            m7  & 51.1           & 2.27        & 0.03              & 130            & 12                   & 0.10 & 0.4         & 9.1        \\
            m8  & 51.1           & 2.18        & 0.02              & 130            & 10                   & 0.09 & 0.4         & 9.1        \\
            m10 & 61.5           & 3.70        & 0.05              & 73             & 3                    & 0.27 & 0.7         & 25.9         \\
            m11 & 61.2           & 3.86        & 0.05              & 80             & 6                    & 0.28 & 0.7         & 27.0        \\
            m12 & 61.2           & 4.27        & 0.06               & 79             & 19                   & 0.33 & 0.8         & 34.1         \\
            m14 & 67.8           & 5.09        & 0.08              & 66             & 6                    & 0.46 & 1.0         & 50.9       \\
            m15 & 67.7           & 5.14        & 0.08               & 56             & 1                    & 0.51 & 1.0         & 50.9         \\
            \hline
        \end{tabular}
    \end{center}

    \caption{Значения вязкостей в экспериментах}
    \label{data}
\end{table}

Как видим, времена и пути релаксации очень малые величины, поэтому предположение
что установившейся скорость достигается на участках 1 и 2 оправдано. Как видим,
числа Рейнольдса в основном меньше 1. Можно предположить что формула Стокса
работает, но окончательный вердикт вынесет график зависимости $ln(\eta)(1/T)$.
Собственно построим график этой зависимости.

Построив график (см. на рис. \ref{graph}) и аппроксимировав точки прямой линией,
получаем следующий угловой коэффициент:
\begin{figure}[H]
    \centering
    \includegraphics{data/graph.pdf}
    \caption{График зависимости $\ln\eta (1 / T)$}
    \label{graph}
\end{figure}

$$W/k = 5.53\cdot10^3\ \text{К} $$

Погрешность будем вычислять как $\sigma_{W/k} = \sqrt{(\sigma_{W/k}^{\text{случ}})^2 +
(\sigma_{W/k}^{\text{сист}})^2}$.

\begin{align*}
\sigma_{W/k}^{\text{случ}} & = \sqrt{\frac{1}{N -
 1}\left(\frac{\left<(\ln\eta)^2\right>}{\left<(1/T)^2\right>}
 - \left(\frac{W}{k}\right)^2 \right)} = 0.09\cdot10^3\ \text{К},\\
\sigma_{W/k}^{\text{сист}} & = \frac{W}{k}\sqrt{\varepsilon_{\ln\eta}^2
        + \varepsilon_{1/T}^2} = 0.15\cdot 10^3\ \text{К}, \\
\frac{W}{k} & = (5.53 \pm 0.17)\cdot 10^3\ \text{К},\ \varepsilon_{W/k} = 3\%,
\end{align*}
В итоге для энергии активации:
$$
W  = (76\pm 2)\cdot 10^{-21}\ \text{Дж},\ \varepsilon_W = 3\%
$$
\section{Обсуждение результатов}

\paragraph{} Проанализируем полученные результаты. На графике
\ref{graph} видно, что значение вязкостей заметно отличаются при низких температурах, при
которых наш метод работает лучше всего. Различие можно объяснить различием
состава глицерина и, возможно, неравномерностью нагрева в нашей установке, т.к.
при низких температурах вязкость меняется резче, и это
может серьезно повлиять на среднюю вязкость.  \end{document}
